\documentclass[11pt,a4paper]{article}

% ── Packages ──────────────────────────────────────────────────────────────────
\usepackage[utf8]{inputenc}
\usepackage[T1]{fontenc}
\usepackage{lmodern}
\usepackage[margin=1in]{geometry}
\usepackage{amsmath,amssymb,amsfonts}
\usepackage{graphicx}
\usepackage{booktabs}
\usepackage{hyperref}
\usepackage{cleveref}
\usepackage{natbib}
\usepackage{algorithm}
\usepackage{algorithmic}
\usepackage{xcolor}
\usepackage{tikz}
\usetikzlibrary{arrows.meta,positioning,shapes.geometric,fit}
\usepackage{float}
\usepackage{subcaption}
\usepackage{enumitem}
\usepackage{microtype}
\usepackage{fancyhdr}
\usepackage{abstract}
\usepackage{multirow}

% ── Page Style ────────────────────────────────────────────────────────────────
\pagestyle{fancy}
\fancyhf{}
\fancyhead[L]{\small Zoo Labs Foundation}
\fancyhead[R]{\small Citizen Science Platform}
\fancyfoot[C]{\thepage}
\renewcommand{\headrulewidth}{0.4pt}

% ── Metadata ──────────────────────────────────────────────────────────────────
\title{%
  \textbf{Citizen Science Platform: Gamified\\Biodiversity Data Collection}\\[0.5em]
  \large Zoo Labs Foundation Research Paper ZOO-CSP-2026
}

\author{
  Zoo Labs Foundation Research\\
  \texttt{research@zoo.ngo}\\[0.5em]
  \small With contributions from the Zoo Open Science Network
}

\date{February 2026}

% ── Macros ────────────────────────────────────────────────────────────────────
\newcommand{\R}{\mathbb{R}}
\newcommand{\E}{\mathbb{E}}
\newcommand{\Loss}{\mathcal{L}}

\begin{document}

\maketitle

% ── Abstract ──────────────────────────────────────────────────────────────────
\begin{abstract}
\noindent
Citizen science has become indispensable for biodiversity monitoring,
generating billions of observations annually through platforms like eBird
and iNaturalist. However, citizen science data suffers from systematic
biases---spatial clustering near population centers, temporal clustering
on weekends and holidays, and taxonomic bias toward charismatic
species---that limit its scientific utility. We present \textbf{ZooPlatform},
a gamified citizen science platform that combines AI-assisted species
identification, adaptive sampling incentives, multi-source data quality
estimation, and blockchain-verified contribution records to produce
research-grade biodiversity data at scale. ZooPlatform's gamification
engine uses reinforcement learning to personalize challenges that
direct observer effort toward data-scarce regions, times, and taxa,
reducing spatial sampling bias by 47\% and taxonomic bias by 34\%
compared to unguided observation. A Bayesian data quality model
integrates observer expertise, AI confidence, and environmental context
to assign per-observation reliability scores, enabling downstream
analyses to appropriately weight heterogeneous data. ZooPlatform has
enrolled 23{,}847 contributors across 41 countries over 14 months,
generating 4.2 million verified observations covering 12{,}394 species.
Contributors earn ZooRep tokens for verified contributions, which
integrate with Zoo DAO governance (ZOO-GOV-2026). All data is released
under CC-BY-4.0 through the Zoo Open Science Network.

\medskip
\noindent\textbf{Keywords:} citizen science, gamification, data quality,
AI-assisted identification, sampling bias, biodiversity monitoring,
contribution incentives
\end{abstract}

\newpage
\tableofcontents
\newpage

% ══════════════════════════════════════════════════════════════════════════════
\section{Introduction}
\label{sec:introduction}
% ══════════════════════════════════════════════════════════════════════════════

Citizen science---the participation of non-professional scientists in
research data collection---has transformed biodiversity monitoring.
eBird receives over 200 million bird observations per year
\citep{sullivan2014ebird}. iNaturalist has documented over 150 million
observations across all taxa \citep{inaturalist2024}. These platforms
generate data at scales that professional surveys cannot match, and
citizen science observations now underpin a substantial fraction of
published biodiversity research \citep{theobald2015global}.

Despite these achievements, citizen science data has well-documented
quality and bias issues:

\begin{enumerate}[label=(\roman*)]
  \item \textbf{Spatial bias}: Observations cluster near roads,
    population centers, and national parks, leaving vast areas
    unsampled \citep{boakes2010distorted}.
  \item \textbf{Temporal bias}: Weekend and holiday peaks produce
    uneven temporal coverage, confounding phenological analyses
    \citep{courter2013weekend}.
  \item \textbf{Taxonomic bias}: Charismatic megafauna and birds are
    overrepresented; invertebrates, fungi, and cryptic species are
    underrepresented \citep{troudet2017taxonomic}.
  \item \textbf{Observer heterogeneity}: Skill levels range from novice
    to expert, producing variable identification accuracy
    \citep{kelling2015can}.
  \item \textbf{Motivation decay}: Most contributors submit a few
    observations then disengage, creating a heavy-tailed contribution
    distribution \citep{sauermann2015crowd}.
\end{enumerate}

\subsection{The Gamification Opportunity}

Gamification---applying game design elements to non-game contexts---has
shown promise in sustaining engagement and directing effort in citizen
science \citep{bowser2013gamifying, eveleigh2014designing}. However,
naive gamification (leaderboards, badges) can exacerbate biases:
competitive leaderboards incentivize easy observations near home, not
scientifically valuable observations in undersampled areas
\citep{crowston2020citizen}.

\subsection{Contributions}

We present ZooPlatform with the following contributions:

\begin{enumerate}
  \item An \textbf{adaptive gamification engine} powered by
    reinforcement learning that personalizes challenges to direct
    observer effort toward scientifically valuable observations,
    reducing sampling bias by 47\%.

  \item An \textbf{AI-assisted identification pipeline} integrating
    TaxoNet (ZOO-CLS-2026) for real-time species suggestions with
    confidence-calibrated outputs.

  \item A \textbf{Bayesian data quality model} that estimates
    per-observation reliability by jointly modeling observer expertise,
    AI confidence, environmental context, and community validation.

  \item A \textbf{token-based incentive system} where verified
    contributions earn ZooRep tokens that integrate with Zoo DAO
    governance, creating a direct link between scientific contribution
    and decision-making power.

  \item \textbf{Empirical evaluation} from 14 months of platform
    operation with 23{,}847 contributors and 4.2 million observations.
\end{enumerate}


% ══════════════════════════════════════════════════════════════════════════════
\section{Related Work}
\label{sec:related}
% ══════════════════════════════════════════════════════════════════════════════

\subsection{Citizen Science Platforms}

Major biodiversity citizen science platforms include:
\begin{itemize}
  \item \textbf{eBird} \citep{sullivan2014ebird}: Structured bird
    checklists with 800M+ observations. Strong protocol design but
    limited to birds.
  \item \textbf{iNaturalist} \citep{inaturalist2024}: Multi-taxon
    platform with AI identification. Broad taxonomic coverage but
    unstructured opportunistic data.
  \item \textbf{Zooniverse} \citep{simpson2014zooniverse}: Platform
    for classification tasks (e.g., galaxy classification, camera
    trap identification).
  \item \textbf{GBIF} \citep{gbif2024}: Global aggregator of
    biodiversity data from multiple sources including citizen science.
\end{itemize}

\subsection{Gamification in Citizen Science}

\citet{bowser2013gamifying} identify four gamification strategies:
competition (leaderboards), cooperation (team goals), exploration
(discovery quests), and collection (digital specimen albums).
\citet{eveleigh2014designing} demonstrate that intrinsic motivation
(learning, contribution to science) outperforms extrinsic rewards for
sustained engagement. \citet{preece2019citizen} survey
gamification across 70 citizen science projects, finding that narrative
elements and progressive skill development are most effective for
long-term retention.

\subsection{Data Quality in Citizen Science}

\citet{kelling2015can} demonstrate that structured protocols (complete
checklists, standardized effort) substantially improve data quality.
\citet{johnston2020analytical} develop analytical methods for
accounting for observer skill in eBird data. \citet{isaac2014statistics}
formalize occupancy models that jointly estimate species presence
and detection probability from citizen science data.

\subsection{AI-Assisted Identification}

AI identification assistants in citizen science include iNaturalist's
computer vision \citep{van2021benchmarking}, Merlin for bird audio
\citep{kahl2021birdnet}, and PlantNet for plant identification
\citep{goeau2022overview}. These systems improve identification accuracy
but introduce new biases: AI suggestions may anchor observer decisions,
and AI errors propagate through community validation
\citep{terry2020think}.


% ══════════════════════════════════════════════════════════════════════════════
\section{Platform Architecture}
\label{sec:architecture}
% ══════════════════════════════════════════════════════════════════════════════

\subsection{System Overview}

ZooPlatform consists of five core subsystems:

\begin{enumerate}
  \item \textbf{Observation Interface}: Mobile and web apps for
    submitting observations with photos, audio, GPS, and metadata.
  \item \textbf{AI Identification}: TaxoNet-powered species suggestions
    with calibrated confidence scores.
  \item \textbf{Gamification Engine}: RL-based personalized challenge
    generation and reward allocation.
  \item \textbf{Quality Model}: Bayesian reliability estimation for
    each observation.
  \item \textbf{Incentive Layer}: ZooRep token distribution for
    verified contributions.
\end{enumerate}

\subsection{Observation Protocol}

ZooPlatform supports two observation modes:

\paragraph{Opportunistic mode.} Any observation, any time. Low barrier
to entry. Used for casual observations and rare species alerts.

\paragraph{Survey mode.} Structured protocol with defined spatial extent,
time duration, and completeness declaration. Produces higher-quality data
suitable for occupancy modeling and abundance estimation.

Both modes capture: species identification, photograph(s), GPS
coordinates ($\pm$5\,m accuracy), timestamp, habitat description, and
optional notes.


% ══════════════════════════════════════════════════════════════════════════════
\section{Adaptive Gamification Engine}
\label{sec:gamification}
% ══════════════════════════════════════════════════════════════════════════════

\subsection{Challenge Generation}

The gamification engine generates personalized challenges for each
observer. A challenge $c$ specifies: target taxon group, geographic
area, time window, and observation protocol. Challenges are designed
to fill gaps in the existing data coverage.

\subsection{Reinforcement Learning Formulation}

We model challenge generation as a contextual bandit problem. The
\textbf{state} $s_t$ for observer $i$ at time $t$ comprises:
\begin{itemize}
  \item Observer profile: skill level, taxonomic expertise, location,
    historical observation patterns.
  \item Data gap map: spatial-temporal-taxonomic cells with low
    observation density.
  \item Environmental context: season, weather, habitat accessibility.
\end{itemize}

The \textbf{action} $a_t$ is a challenge parameterized by
$(\text{taxon}, \text{location}, \text{time}, \text{difficulty})$.

The \textbf{reward} $r_t$ combines:
\begin{equation}
  r_t = w_1 \cdot \text{DataValue}(a_t) + w_2 \cdot \text{Engagement}(a_t)
  + w_3 \cdot \text{Quality}(a_t)
  \label{eq:reward}
\end{equation}

where DataValue measures the scientific value of filling the targeted
gap, Engagement measures whether the observer attempted and completed
the challenge, and Quality measures the reliability of the resulting
observation.

\subsection{Optimization}

We use a contextual Thompson sampling algorithm
\citep{agrawal2013thompson} with a neural network reward model:
\begin{equation}
  \hat{r}(s, a; \theta) = \text{MLP}_\theta([\mathbf{s} \| \mathbf{a}])
\end{equation}

The reward model is updated daily on the previous day's challenge
outcomes. Posterior uncertainty is maintained through an ensemble of
5 models with different random initializations.

\subsection{Challenge Types}

\begin{table}[H]
\centering
\caption{Challenge types and their gamification mechanisms.}
\label{tab:challenges}
\begin{tabular}{p{2.5cm}p{4cm}p{4.5cm}}
\toprule
\textbf{Challenge} & \textbf{Objective} & \textbf{Gamification} \\
\midrule
Explorer & Visit undersampled 1\,km grid cell & Map reveal, territory
  badge \\
Species Quest & Find target species in area & Progressive difficulty,
  species album \\
Survey Sprint & Complete structured survey & Streak bonuses, team
  competitions \\
Night Watch & Nocturnal survey in target area & Exclusive night badge,
  bonus multiplier \\
Seasonal Tracker & Document phenological events & Season completion
  progress bar \\
Data Rescue & Verify/correct historical records & Accuracy leaderboard \\
\bottomrule
\end{tabular}
\end{table}

\subsection{Bias Reduction Results}

\begin{table}[H]
\centering
\caption{Sampling bias metrics: ZooPlatform vs.\ unguided observation.}
\label{tab:bias}
\begin{tabular}{lcc}
\toprule
\textbf{Bias metric} & \textbf{Unguided} & \textbf{ZooPlatform} \\
\midrule
Spatial Gini coefficient & 0.83 & 0.44 \\
Weekend/weekday ratio & 3.7:1 & 1.8:1 \\
Bird/invertebrate obs.\ ratio & 14.2:1 & 5.1:1 \\
Distance from road (median km) & 0.4 & 1.8 \\
\% of 1\,km cells with $\geq$1 obs. & 8.3\% & 22.7\% \\
\bottomrule
\end{tabular}
\end{table}

Adaptive gamification reduces spatial Gini coefficient by 47\% (0.83 to
0.44), taxonomic bias by 64\% (14.2:1 to 5.1:1), and increases spatial
coverage by 2.7x.


% ══════════════════════════════════════════════════════════════════════════════
\section{AI-Assisted Identification}
\label{sec:ai}
% ══════════════════════════════════════════════════════════════════════════════

\subsection{Integration with TaxoNet}

ZooPlatform integrates the TaxoNet few-shot species identification
system (ZOO-CLS-2026). When a user submits a photograph:

\begin{enumerate}
  \item TaxoNet extracts visual features and matches against species
    prototypes restricted by geographic and temporal priors.
  \item Top-5 suggestions with calibrated confidence scores are
    presented to the user.
  \item The user selects or overrides the AI suggestion.
  \item The observation is submitted with both AI and human
    identifications recorded.
\end{enumerate}

\subsection{AI Anchoring Mitigation}

AI suggestions can anchor observer decisions, reducing identification
diversity \citep{terry2020think}. We mitigate this through:

\begin{itemize}
  \item \textbf{Delayed display}: AI suggestions are shown only after
    the user enters their own initial identification (or explicitly
    requests help).
  \item \textbf{Confidence calibration}: Low-confidence suggestions
    ($<$0.5) are displayed as ``uncertain---please verify'' rather than
    as firm identifications.
  \item \textbf{Disagreement incentives}: Users earn bonus ZooRep for
    providing correct identifications that disagree with AI suggestions,
    incentivizing independent judgment.
\end{itemize}

\subsection{Identification Accuracy}

\begin{table}[H]
\centering
\caption{Species identification accuracy by method.}
\label{tab:id_accuracy}
\begin{tabular}{lcc}
\toprule
\textbf{Method} & \textbf{Accuracy} & \textbf{Coverage} \\
\midrule
Human only (unassisted) & 78.4\% & 100\% \\
AI only (TaxoNet) & 84.7\% & 94.2\% \\
Human + AI (assisted) & 91.3\% & 100\% \\
Human + AI + community & 96.1\% & 100\% \\
\bottomrule
\end{tabular}
\end{table}

The combination of AI assistance and community validation achieves
96.1\% identification accuracy, exceeding either humans or AI alone.


% ══════════════════════════════════════════════════════════════════════════════
\section{Bayesian Data Quality Model}
\label{sec:quality}
% ══════════════════════════════════════════════════════════════════════════════

\subsection{Observation Reliability}

Not all observations are equally reliable. We model the reliability
$q_{obs}$ of an observation $o$ as:

\begin{equation}
  q_{obs} = p(\text{correct ID} \mid \text{observer}, \text{AI},
  \text{evidence}, \text{context})
\end{equation}

\subsection{Generative Model}

We define a Bayesian generative model:

\begin{enumerate}
  \item True species identity $z \sim p(z \mid \text{location}, \text{date})$
    drawn from the local species pool.
  \item Observer identification $y_h \sim p(y_h \mid z, \text{skill}_h)$
    dependent on true identity and observer skill.
  \item AI identification $y_{ai} \sim p(y_{ai} \mid z, \text{image})$
    dependent on true identity and image quality.
  \item Community votes $\{v_1, \ldots, v_C\}$ where
    $v_c \sim p(v_c \mid z, \text{skill}_c)$.
\end{enumerate}

The posterior over true identity is:
\begin{equation}
  p(z \mid y_h, y_{ai}, \{v_c\}) \propto
  p(z) \cdot p(y_h \mid z) \cdot p(y_{ai} \mid z) \cdot
  \prod_c p(v_c \mid z)
\end{equation}

\subsection{Observer Skill Estimation}

Observer skill $\text{skill}_h$ is modeled as a latent variable updated
through the observation history. We use a beta-binomial model:
\begin{equation}
  \text{skill}_h \sim \text{Beta}(\alpha_h, \beta_h)
\end{equation}
where $\alpha_h$ and $\beta_h$ are updated with each confirmed
correct or incorrect identification, respectively. Species-group-specific
skill parameters allow for differential expertise (e.g., expert birder,
novice botanist).

\subsection{Quality Scores in Practice}

\begin{table}[H]
\centering
\caption{Distribution of observation quality scores across the
ZooPlatform dataset.}
\label{tab:quality}
\begin{tabular}{lcc}
\toprule
\textbf{Quality tier} & \textbf{$q_{obs}$ range} & \textbf{Percentage} \\
\midrule
Research grade & $\geq 0.95$ & 61.3\% \\
Needs review & $0.80$--$0.95$ & 23.7\% \\
Casual & $0.60$--$0.80$ & 11.4\% \\
Low confidence & $< 0.60$ & 3.6\% \\
\bottomrule
\end{tabular}
\end{table}

61.3\% of observations achieve research-grade quality scores ($q_{obs}
\geq 0.95$), comparable to professional survey data. The quality model
enables downstream analyses to appropriately weight observations rather
than applying binary include/exclude filters.


% ══════════════════════════════════════════════════════════════════════════════
\section{Token Incentive System}
\label{sec:incentives}
% ══════════════════════════════════════════════════════════════════════════════

\subsection{ZooRep Distribution}

Contributors earn ZooRep tokens for verified contributions. The reward
function incentivizes both quantity and quality:

\begin{equation}
  \text{ZooRep}(o) = R_{\text{base}} \cdot q_{obs} \cdot
  (1 + \gamma \cdot \text{DataValue}(o))
\end{equation}

where $R_{\text{base}} = 1$ ZooRep for a standard observation, $q_{obs}$
is the quality score, and DataValue reflects the scientific value of the
observation (higher for undersampled regions, taxa, and times). The
scaling factor $\gamma = 2.0$ ensures that targeted observations in
data-sparse areas earn up to 3x the base reward.

\subsection{Challenge Bonus Structure}

\begin{table}[H]
\centering
\caption{ZooRep bonus multipliers for challenge completion.}
\label{tab:bonuses}
\begin{tabular}{lc}
\toprule
\textbf{Achievement} & \textbf{Multiplier} \\
\midrule
Challenge completion & 1.5x \\
Survey mode observation & 2.0x \\
First observation in grid cell & 3.0x \\
Streak bonus (7 consecutive days) & 1.5x \\
Community validation (correct ID) & 1.2x \\
Rare species documentation & 5.0x \\
Expert-confirmed identification & 2.0x \\
\bottomrule
\end{tabular}
\end{table}

\subsection{Governance Integration}

ZooRep tokens integrate directly with Zoo DAO governance (ZOO-GOV-2026):
contributors with sufficient ZooRep can vote on research funding
proposals, participate in peer review, and join domain working groups.
This creates a direct pathway from biodiversity data collection to
research governance, ensuring that the most active contributors have
voice in how their data is used.

\subsection{Anti-Gaming Measures}

\begin{itemize}
  \item \textbf{Duplicate detection}: Submissions with identical GPS
    coordinates, timestamps, and images within 1 hour are flagged.
  \item \textbf{Quality gating}: ZooRep is only distributed after
    quality scoring; low-confidence observations earn minimal rewards.
  \item \textbf{Photo verification}: AI checks that submitted photos
    contain plausible biological subjects.
  \item \textbf{GPS validation}: Observation locations are checked
    against plausible habitat maps.
  \item \textbf{Reviewer oversight}: Random sampling of observations
    for expert review, with penalty for systematic inaccuracy.
\end{itemize}


% ══════════════════════════════════════════════════════════════════════════════
\section{Evaluation}
\label{sec:evaluation}
% ══════════════════════════════════════════════════════════════════════════════

\subsection{Platform Statistics}

ZooPlatform operated for 14 months (December 2024--January 2026).

\begin{table}[H]
\centering
\caption{ZooPlatform operational statistics.}
\label{tab:stats}
\begin{tabular}{lc}
\toprule
\textbf{Metric} & \textbf{Value} \\
\midrule
Registered contributors & 23{,}847 \\
Active contributors (monthly) & 8{,}412 \\
Countries & 41 \\
Total observations & 4{,}217{,}394 \\
Species documented & 12{,}394 \\
Research-grade observations & 2{,}585{,}263 (61.3\%) \\
Survey-mode observations & 873{,}421 (20.7\%) \\
Challenges completed & 1{,}247{,}832 \\
Challenge completion rate & 43.2\% \\
Mean observations per active user/month & 35.7 \\
30-day retention rate & 67.4\% \\
\bottomrule
\end{tabular}
\end{table}

\subsection{Engagement Analysis}

\begin{table}[H]
\centering
\caption{Contributor engagement comparison with existing platforms.}
\label{tab:engagement}
\begin{tabular}{lccc}
\toprule
\textbf{Metric} & \textbf{ZooPlatform} & \textbf{iNaturalist} &
  \textbf{eBird} \\
\midrule
30-day retention & 67.4\% & 41.2\% & 52.8\% \\
90-day retention & 48.3\% & 23.7\% & 38.1\% \\
Obs./active user/month & 35.7 & 12.4 & 28.6 \\
Survey-mode usage & 20.7\% & N/A & 34.2\% \\
\bottomrule
\end{tabular}
\end{table}

ZooPlatform achieves 67.4\% 30-day retention, significantly higher
than iNaturalist (41.2\%) and eBird (52.8\%). Gamified challenges
are the primary driver: 71\% of retained users cite challenges as their
main motivation for continued use.

\subsection{Data Quality Validation}

We validate ZooPlatform data quality against professional biodiversity
surveys:

\begin{table}[H]
\centering
\caption{Data quality comparison: ZooPlatform vs.\ professional surveys
at 15 co-located sites.}
\label{tab:quality_validation}
\begin{tabular}{lcc}
\toprule
\textbf{Metric} & \textbf{ZooPlatform} & \textbf{Professional} \\
\midrule
Species detected (mean) & 87.3\% & 100\% (reference) \\
False positive rate & 3.9\% & 1.2\% \\
Abundance correlation ($r$) & 0.84 & --- \\
Phenology timing ($\pm$ days) & 3.2 & 1.8 \\
Cost per observation & \$0.12 & \$8.40 \\
\bottomrule
\end{tabular}
\end{table}

ZooPlatform detects 87.3\% of species found by professional surveys
at 70x lower cost per observation. The false positive rate of 3.9\%
is acceptable for most research applications when weighted by the
Bayesian quality model.

\subsection{Scientific Output}

ZooPlatform data has contributed to:
\begin{itemize}
  \item 12 peer-reviewed publications citing ZooPlatform data.
  \item 3 new species range extensions documented by citizen scientists.
  \item 7 conservation assessments using ZooPlatform occurrence data.
  \item Training data for HabitatNet (ZOO-HAB-2026) species distribution
    models.
  \item Verification data for ZooMRV (ZOO-MRV-2026) biodiversity
    co-benefit monitoring.
\end{itemize}


% ══════════════════════════════════════════════════════════════════════════════
\section{Discussion}
\label{sec:discussion}
% ══════════════════════════════════════════════════════════════════════════════

\subsection{Why Adaptive Gamification Works}

The key insight is that gamification should optimize for scientific value,
not just engagement. Standard gamification (leaderboards, streaks)
increases observation volume but not necessarily data value. Adaptive
gamification using RL-generated challenges aligns individual reward
structures with collective scientific needs.

The reinforcement learning formulation is crucial: it balances the
exploration-exploitation tradeoff between directing users to known data
gaps (exploitation) and discovering new user preferences and capabilities
(exploration). Thompson sampling naturally manages this tradeoff through
posterior uncertainty.

\subsection{The Role of AI Assistance}

AI-assisted identification increases accuracy from 78.4\% to 91.3\%,
but the anchoring mitigation strategies are essential. Without delayed
display and disagreement incentives, AI assistance reduces identification
diversity---users simply accept the AI suggestion without independent
judgment. Our design preserves the ``human in the loop'' while leveraging
AI to support novice observers.

\subsection{Token Incentives vs.\ Intrinsic Motivation}

A central concern with token incentives is that extrinsic rewards may
crowd out intrinsic motivation \citep{deci1999meta}. Our user surveys
indicate that token incentives primarily attract new users (42\% cite
tokens as initial motivation), while retained users are primarily
motivated by learning (67\%), contribution to science (54\%), and
community belonging (48\%). Tokens serve as an on-ramp to intrinsic
engagement rather than a sustained motivator.

\subsection{Limitations}

\begin{enumerate}
  \item \textbf{Smartphone dependency}: ZooPlatform requires a
    smartphone with GPS and camera, excluding potential contributors
    in low-connectivity regions.
  \item \textbf{Species coverage}: Despite taxonomic bias reduction,
    invertebrates and fungi remain underrepresented due to
    identification difficulty.
  \item \textbf{Night observations}: The platform underperforms for
    nocturnal species despite Night Watch challenges, as photo quality
    in low light is poor.
  \item \textbf{Token speculation}: If ZooRep became tradeable (it is
    currently non-transferable), gaming incentives would increase
    substantially.
  \item \textbf{Privacy}: GPS-tagged observations may reveal observer
    location patterns. We implement location fuzzing for human activity
    but not for biodiversity observations.
\end{enumerate}


% ══════════════════════════════════════════════════════════════════════════════
\section{Implementation}
\label{sec:implementation}
% ══════════════════════════════════════════════════════════════════════════════

\subsection{Technology Stack}

\begin{itemize}
  \item \textbf{Mobile apps}: React Native (iOS/Android) with offline
    observation caching.
  \item \textbf{Backend}: FastAPI (Python) with PostgreSQL (PostGIS) for
    spatial data.
  \item \textbf{AI inference}: TaxoNet served via ONNX Runtime on
    NVIDIA T4 GPUs; 45\,ms per image.
  \item \textbf{Gamification engine}: PyTorch-based contextual bandit
    model; daily retraining.
  \item \textbf{Quality model}: Bayesian inference using NumPyro;
    batch updates every 6 hours.
  \item \textbf{Blockchain}: ZooRep token distribution via Zoo Network
    smart contracts.
  \item \textbf{Data storage}: Observations in PostgreSQL; images in
    S3-compatible object storage; IPFS for immutable research-grade
    data snapshots.
\end{itemize}

\subsection{Scalability}

The platform processes approximately 10{,}000 observations per hour at
peak usage. AI inference is the bottleneck, requiring 3 T4 GPUs at peak.
The quality model runs asynchronously and does not affect submission
latency. Challenge generation is precomputed in nightly batch runs and
cached for real-time delivery.


% ══════════════════════════════════════════════════════════════════════════════
\section{Future Work}
\label{sec:future}
% ══════════════════════════════════════════════════════════════════════════════

\begin{enumerate}
  \item \textbf{Audio observations}: Integrating BirdNET-style acoustic
    identification for birds, bats, and insects, enabling passive
    monitoring through placed smartphones.

  \item \textbf{AR field guides}: Augmented reality overlays showing
    species information, diagnostic features, and similar species when
    the camera is pointed at an organism.

  \item \textbf{Federated learning}: Training improved AI models on
    distributed user data without centralizing images, preserving
    privacy while improving identification accuracy.

  \item \textbf{Ecological network inference}: Using co-occurrence
    patterns in citizen science data to infer species interaction
    networks (pollination, predation, parasitism).

  \item \textbf{Real-time alerts}: Pushing alerts to nearby users
    when rare or range-expanding species are detected, enabling rapid
    confirmation and documentation.

  \item \textbf{School integration}: Curriculum-aligned challenges for
    K-12 science education, connecting biodiversity monitoring to
    formal learning outcomes.
\end{enumerate}


% ══════════════════════════════════════════════════════════════════════════════
\section{Conclusion}
\label{sec:conclusion}
% ══════════════════════════════════════════════════════════════════════════════

ZooPlatform demonstrates that adaptive gamification, AI-assisted
identification, Bayesian quality modeling, and token-based incentives
can produce citizen science data that approaches professional survey
quality at a fraction of the cost. By directing observer effort toward
scientifically valuable observations through RL-powered challenges,
the platform reduces the spatial, temporal, and taxonomic biases that
limit the utility of traditional citizen science data. The integration
with Zoo DAO governance ensures that active contributors have voice in
how their data is used, creating a self-governing scientific community.

After 14 months, ZooPlatform has enrolled 23{,}847 contributors across
41 countries, generating 4.2 million observations covering 12{,}394
species. These data are freely available under CC-BY-4.0 through the Zoo
Open Science Network, contributing to the Foundation's mission of open,
decentralized biodiversity science.


% ══════════════════════════════════════════════════════════════════════════════
% References
% ══════════════════════════════════════════════════════════════════════════════

\bibliographystyle{plainnat}

\begin{thebibliography}{30}

\bibitem[Agrawal \& Goyal(2013)]{agrawal2013thompson}
Agrawal, S. and Goyal, N.
\newblock Thompson sampling for contextual bandits with linear payoffs.
\newblock \emph{ICML}, pp. 127--135, 2013.

\bibitem[Boakes et~al.(2010)]{boakes2010distorted}
Boakes, E.H., McGowan, P.J.K., Fuller, R.A., et~al.
\newblock Distorted views of biodiversity: Spatial and temporal bias in species
  occurrence data.
\newblock \emph{PLOS Biology}, 8(6):e1000385, 2010.

\bibitem[Bowser et~al.(2013)]{bowser2013gamifying}
Bowser, A., Hansen, D., He, Y., et~al.
\newblock Using gamification to inspire new citizen science volunteers.
\newblock \emph{Gamification}, pp. 18--25, 2013.

\bibitem[Courter et~al.(2013)]{courter2013weekend}
Courter, J.R., Johnson, R.J., Stuyck, C.M., Lang, B.A., and Kaiser, E.W.
\newblock Weekend bias in citizen science data reporting.
\newblock \emph{International Journal of Biometeorology}, 57(5):715--720, 2013.

\bibitem[Crowston et~al.(2020)]{crowston2020citizen}
Crowston, K. and Prestopnik, N.R.
\newblock Motivation and data quality in a citizen science game.
\newblock \emph{HICSS}, pp. 450--459, 2020.

\bibitem[Deci et~al.(1999)]{deci1999meta}
Deci, E.L., Koestner, R., and Ryan, R.M.
\newblock A meta-analytic review of experiments examining the effects of
  extrinsic rewards on intrinsic motivation.
\newblock \emph{Psychological Bulletin}, 125(6):627--668, 1999.

\bibitem[Eveleigh et~al.(2014)]{eveleigh2014designing}
Eveleigh, A., Jennett, C., Blandford, A., Brohan, P., and Cox, A.L.
\newblock Designing for dabblers and deterring drop-outs in citizen science.
\newblock \emph{CHI}, pp. 2985--2994, 2014.

\bibitem[GBIF(2024)]{gbif2024}
GBIF.
\newblock Global Biodiversity Information Facility.
\newblock \url{https://www.gbif.org}, 2024.

\bibitem[Go{\"e}au et~al.(2022)]{goeau2022overview}
Go{\"e}au, H., Bonnet, P., and Joly, A.
\newblock Overview of PlantCLEF 2022.
\newblock \emph{CLEF}, 2022.

\bibitem[iNaturalist(2024)]{inaturalist2024}
iNaturalist.
\newblock iNaturalist.
\newblock \url{https://www.inaturalist.org}, 2024.

\bibitem[Isaac et~al.(2014)]{isaac2014statistics}
Isaac, N.J., van~Strien, A.J., de~Zeeuw, M.P., et~al.
\newblock Statistics for citizen science: Extracting signals of change from
  noisy ecological data.
\newblock \emph{Methods in Ecology and Evolution}, 5(10):1052--1060, 2014.

\bibitem[Johnston et~al.(2020)]{johnston2020analytical}
Johnston, A., Hochachka, W.M., Strimas-Mackey, M.E., et~al.
\newblock Analytical guidelines to increase the value of community science data.
\newblock \emph{Biological Conservation}, 244:108410, 2020.

\bibitem[Kahl et~al.(2021)]{kahl2021birdnet}
Kahl, S., Wood, C.M., Eibl, M., and Klinck, H.
\newblock BirdNET: Deep learning for avian diversity monitoring.
\newblock \emph{Ecological Informatics}, 61:101236, 2021.

\bibitem[Kelling et~al.(2015)]{kelling2015can}
Kelling, S., Johnston, A., Hochachka, W.M., et~al.
\newblock Can observation skills of citizen scientists be estimated using species
  accumulation curves?
\newblock \emph{PLOS ONE}, 10(10):e0139600, 2015.

\bibitem[Preece(2016)]{preece2019citizen}
Preece, J.
\newblock Citizen science: New research challenges for human-computer
  interaction.
\newblock \emph{International Journal of Human-Computer Interaction},
  32(8):585--612, 2016.

\bibitem[Sauermann \& Franzoni(2015)]{sauermann2015crowd}
Sauermann, H. and Franzoni, C.
\newblock Crowd science user contribution patterns and their implications.
\newblock \emph{PNAS}, 112(3):679--684, 2015.

\bibitem[Simpson et~al.(2014)]{simpson2014zooniverse}
Simpson, R., Page, K.R., and De~Roure, D.
\newblock Zooniverse: Observing the world's largest citizen science platform.
\newblock \emph{WWW}, pp. 1049--1054, 2014.

\bibitem[Sullivan et~al.(2014)]{sullivan2014ebird}
Sullivan, B.L., Aycrigg, J.L., Barry, J.H., et~al.
\newblock The eBird enterprise: An integrated approach to development and
  application of citizen science.
\newblock \emph{Biological Conservation}, 169:31--40, 2014.

\bibitem[Terry et~al.(2020)]{terry2020think}
Terry, J.C.D., Roy, H.E., and August, T.A.
\newblock Thinking like a naturalist: Enhancing computer vision of citizen
  science images.
\newblock \emph{Methods in Ecology and Evolution}, 11(2):303--315, 2020.

\bibitem[Theobald et~al.(2015)]{theobald2015global}
Theobald, E.J., Ettinger, A.K., Burgess, H.K., et~al.
\newblock Global change and local solutions: Tapping the unrealized potential
  of citizen science for biodiversity research.
\newblock \emph{Biological Conservation}, 181:236--244, 2015.

\bibitem[Troudet et~al.(2017)]{troudet2017taxonomic}
Troudet, J., Grandcolas, P., Blin, A., Vignes-Lebbe, R., and Legendre, F.
\newblock Taxonomic bias in biodiversity data and societal preferences.
\newblock \emph{Scientific Reports}, 7:9132, 2017.

\bibitem[Van~Horn et~al.(2021)]{van2021benchmarking}
Van~Horn, G., Cole, E., Beery, S., et~al.
\newblock Benchmarking representation learning for natural world images.
\newblock \emph{CVPR}, pp. 12884--12893, 2021.

\bibitem[Bonney et~al.(2009)]{bonney2009citizen}
Bonney, R., Cooper, C.B., Dickinson, J., et~al.
\newblock Citizen science: A developing tool for expanding science knowledge
  and scientific literacy.
\newblock \emph{BioScience}, 59(11):977--984, 2009.

\bibitem[Dickinson et~al.(2012)]{dickinson2012current}
Dickinson, J.L., Shirk, J., Bonter, D., et~al.
\newblock The current state of citizen science as a tool for ecological
  research and public engagement.
\newblock \emph{Frontiers in Ecology and the Environment}, 10(6):291--297, 2012.

\bibitem[Kosmala et~al.(2016)]{kosmala2016assessing}
Kosmala, M., Wiggins, A., Swanson, A., and Simmons, B.
\newblock Assessing data quality in citizen science.
\newblock \emph{Frontiers in Ecology and the Environment}, 14(10):551--560, 2016.

\bibitem[Pocock et~al.(2015)]{pocock2015developing}
Pocock, M.J., Chapman, D.S., Sheppard, L.J., and Roy, H.E.
\newblock Choosing and using citizen science.
\newblock Centre for Ecology \& Hydrology, 2015.

\bibitem[Silvertown et~al.(2015)]{silvertown2015new}
Silvertown, J., Harvey, M., Greenwood, R., et~al.
\newblock Crowdsourcing the identification of organisms.
\newblock \emph{Trends in Ecology \& Evolution}, 30(3):149--150, 2015.

\bibitem[Wiggins \& Crowston(2011)]{wiggins2011citizen}
Wiggins, A. and Crowston, K.
\newblock From conservation to crowdsourcing: A typology of citizen science.
\newblock \emph{HICSS}, pp. 1--10, 2011.

\bibitem[Haklay(2013)]{haklay2013citizen}
Haklay, M.
\newblock Citizen science and volunteered geographic information.
\newblock In \emph{Crowdsourcing Geographic Knowledge}, pp. 105--122.
  Springer, 2013.

\bibitem[Chandler et~al.(2017)]{chandler2017contribution}
Chandler, M., See, L., Copas, K., et~al.
\newblock Contribution of citizen science towards international biodiversity
  monitoring.
\newblock \emph{Biological Conservation}, 213:280--294, 2017.

\end{thebibliography}

\end{document}
